% §2.5 Bibliographic Notes
%
% Cross-reference resolution status:
%
% ALREADY RESOLVED (within chapter 2, drafted in Tasks 2-5 of Plan 2):
%   \Cref{sec:bsc}                  -- chapter 2 §2.1
%   \Cref{sec:asymmetric-channel}   -- chapter 2 §2.2
%   \Cref{sec:channel-composition}  -- chapter 2 §2.3
%   \Cref{sec:model-order}          -- chapter 2 §2.4
%   \Cref{def:model-order}          -- chapter 2 §2.4 Definition 2.x
%   \Cref{thm:boolean-ceiling}      -- chapter 2 §2.4 Boolean Ceiling Theorem
%
% ALREADY RESOLVED (chapter 4, drafted in Plan 1):
%   \Cref{sec:bib-notes-ch4}        -- chapter 4 §4.7 Bibliographic Notes

\section*{Bibliographic Notes}
\addcontentsline{toc}{section}{Bibliographic Notes}
\label{sec:bib-notes-ch2}

\paragraph{Shannon and the channel framework.}
The channel matrix formalism used throughout this chapter descends directly
from \textcite{shannon-1948}. Shannon's paper introduced the discrete
memoryless channel as the mathematical object for studying noisy
communication: a channel is specified by a conditional distribution
$P(\text{output} \mid \text{input})$, and the capacity of the channel is the
maximum mutual information between input and output over all input
distributions. The binary symmetric channel of \Cref{sec:bsc} is one of
Shannon's central running examples; the crossover parameter $\fprate$ and
the capacity formula $1 - H(\fprate)$ are both in that 1948 paper. The
treatment of the asymmetric binary channel in \Cref{sec:asymmetric-channel}
is a natural extension of Shannon's framework to the case where the two
input symbols face independent error probabilities. For a modern textbook
development of binary channels, channel capacity, and the channel coding
theorem, see \textcite{cover-thomas}, especially Chapters~7 and~8.

\paragraph{The Bernoulli papers.}
The connection between binary channels and the Bernoulli framework is drawn
in \textcite{towell-bernoulli-sets}, which introduces the two-axiom model
and proves the composition theorem that \Cref{sec:channel-composition}
reformulates in matrix-product language. The matrix-product form, that
serial composition of two binary channels corresponds to matrix multiplication
of their channel matrices, is a reorganization of the rate-propagation
formula in that paper; see the companion notes in \Cref{sec:bib-notes-ch4}
for a precise account of what is new in this text and what descends from the
companion papers.

The model-order treatment of \Cref{sec:model-order}, and in particular the
definition of model order as the number of partition blocks in the coarsest
compatible partition of the input domain (\Cref{def:model-order}), follows
\textcite{towell-bernoulli-data-type}. That paper extends the Bernoulli
framework from Boolean sets to approximate values over arbitrary algebraic
types; the order-2 Boolean case developed in this chapter is the simplest
instance of its general theory. The Boolean Ceiling Theorem
(\Cref{thm:boolean-ceiling}), that any binary channel over $\{\top, \bot\}$
has model order at most $2$, is explicit in the algebraic-types paper as
the base case of the more general statement that a channel over a $k$-element
domain has model order at most $k$.

\paragraph{Model order versus degrees of freedom.}
A clarification that deserves explicit attention is the distinction between
model order and parameter count, developed in \Cref{sec:model-order} under
the heading ``Model order vs.\ parameter count.'' This distinction is easy
to miss in the source papers, because for a single Boolean query the two
quantities agree: an $n$-th order channel on one Boolean input has $n$
partition blocks and also $n$ free parameters, one per block. The difference
surfaces in the multi-query setting. When $m$ queries
are issued jointly, a general order-$n$ channel on $m$ inputs has
$n(2^m - 1)$ parameters in the per-source-block convention, but the
independence axioms (Axioms~1 and~2) force a Kronecker factorization that
reduces this to $nm$. The model order does not change: it remains $n$,
determined by the partition of the single-query input domain. What changes
is the parameter count, and the axioms are entirely responsible for the
reduction. Model order constrains which rows of the channel matrix $Q$ can
be identified; the axioms constrain how the channel matrix factorizes across
queries. These are orthogonal constraints.

\paragraph{Further reading.}
Chapter~4's Bibliographic Notes (\Cref{sec:bib-notes-ch4}) provide
additional context for the Kronecker product and its role in the multi-query
factorization, and for the relationship between the Bernoulli axioms and the
memoryless channel. Readers who want a deeper grounding in information
theory beyond what Part~I requires will find \textcite{cover-thomas}
the standard reference.
